\documentclass[book.tex]{subfiles}
\begin{document}

\chapter{Inverse Design}

\label{chap:selfies}
\noindent\rule{11cm}{0.4pt} \\
\textbf{Prerequisites:} \autoref{chap:graphconvs}, \autoref{chap:equivariant_networks}, \autoref{chap:molecular_dynamics},  \autoref{chap:dft} \\
\textbf{Difficulty Level:} ** \\
\noindent\rule{11cm}{0.4pt}
\vspace{5mm}

"Deep learning has proven to yield fast and accurate predictions of quantum-chemical properties to accelerate the discovery of novel molecules and materials. As an exhaustive exploration of the vast chemical space is still infeasible, we require generative models that guide our search towards systems with desired properties. Many graph-based models have previously been proposed to tackle this challenge. However, they are restricted by a lack of spatial information such that they are unable to recognize spatial isomerism and non-bonded interactions. Generative models for 3-dimensional structures are an emerging field promising to overcome these restrictions. In this talk we will introduce the challenges of 3d molecule generation and focus on the conditional autoregressive deep neural network cG-SchNet, which generates molecules with specified chemical and structural properties by placing one atom after another in 3d space. The model enables targeted sampling of novel structures from conditional distributions, even in domains where reference calculations are sparse."

References \cite{gebauer2022inverse}


Notes: Have equivariant methods and flows. Have one shot methods.

Proposes an autoregressive structure based on gSchNet. This work adds conditions such as structural constraints. onto the base architecture. Architecture is cGSchNet. SchNet is rotation, permutation, and translation invariant. Do a Gaussian expansion. Have a condition embedding network

\begin{align}
(\lambda_1,\dotsc,\lambda_k) \mapsto y
\end{align}

Have an atom type prediction network. We have an atom placement loop.

\begin{enumerate}
\item Randomly choose an atom as focus from atoms which aren't marked as finished.
\item Predict and sample the type $Z_{next}$ of the next atom. Include a stop type to cut off the generation.
\item Predict the probabilities of pairwise distribution of distances between next atom and preceding atoms or tokens. This is done with a discretized/binned Gaussian?
\item Look up the probabilities of the pairwise distances.
\item Multiply probabilities at each position.
\end{enumerate}

There are some approximation artifacts here.

One limitation is models are primarily only on QM9 and other small datasets.

What is the origin token?

GSchNet is generative SchNet.
\end{document}